Predicting properties of trained neural networks from their weights --- model zoo analysis ---
typically assumes that test models share the same architecture as training models.
When this assumption breaks, flat tokenization methods collapse:
we show that SANE, the current state of the art for within-architecture weight-space SSL,
maps all ViT inputs to near-identical latent codes after training on CNN checkpoints,
achieving $R^2 = 0.072$ (near-chance) on 53 held-out ViT-S/16 ImageNet checkpoints.
We propose that the root cause is the \emph{representation substrate}, not the learning algorithm:
flat weight tokenization discards the relational graph structure that is consistent across architecture families.
We train a permutation-equivariant graph SSL encoder on MLP/CNN model zoos
and demonstrate zero-shot cross-architecture transfer to ViT property prediction,
achieving $R^2 = 0.231$ versus SANE's $R^2 = 0.072$ --- a $3\times$ improvement ---
without any ViT training data.
Counterintuitively, adding scale equivariance to the symmetry group hurts rather than helps ---
we attribute this to a \emph{LayerNorm gauge-fixing hypothesis}:
LayerNorm eliminates the scale degree of freedom that scale equivariance is designed to handle,
making scale invariance an incorrect inductive bias for ViT weight encoding.
These results suggest that the coordinate system (computational graph schema)
matters more than the symmetry group for cross-architecture weight representation learning.
